\documentclass[11pt]{article}
\usepackage[margin=1in]{geometry}
\usepackage{graphicx}
\usepackage{booktabs}
\usepackage{array}
\usepackage{hyperref}
\usepackage{float}

\title{MNIST BNN Power Trace Capture on CW305 + CW-Lite\\
\large Output Summary, Transformation Evidence, and Failure Review}
\author{Hardware evaluation notes}
\date{July 5, 2026}

\begin{document}
\maketitle

\begin{abstract}
This report summarizes the corrected MNIST hardware run for the BNN side-channel
reconstruction experiment.  The target is MNIST only, not CIFAR.  The FPGA was
programmed with the CW305 Artix-7 bitstream, the CW-Lite captured analog power
traces, and the first convolution output write was used as the trigger anchor.
The functional hardware result is correct: the CW305 output matches the software
golden model bit-for-bit for all 784 MNIST pixels.  The attack result is still
poor: the classifier recognizes original images at 97.5\%, but recognizes the
curve-fit background reconstructions at only 8.5\% and curve-fit template
reconstructions at 13.0\%.  The bottleneck is therefore the trace-to-image
recovery stage, not MNIST classification or the basic HLS functional behavior.
\end{abstract}

\section{Scope}
The present run uses the MNIST data path and the MNIST layer-1 BNN parameters.
CIFAR files and CIFAR shapes are deliberately excluded.  The relevant hardware
configuration is:

\begin{itemize}
  \item Target board: CW305 Artix-7.
  \item Capture board: ChipWhisperer-Lite.
  \item Bitstream: \path{build/cw305_bnn_3_f16_trigfix.bit}.
  \item Captures: \path{host/traces_analog_trigfix_500}.
  \item Evaluation without curve fitting: \path{attack/results_trigfix500_nofit_norm}.
  \item Evaluation with curve fitting: \path{attack/results_trigfix500_curvefit_norm}.
\end{itemize}

\section{Hardware and Capture Summary}
The capture procedure was changed from the older busy-window trigger to a trigger
on the first HLS output write.  This matters because the HLS design spends a
large number of cycles in initialization and control before useful MNIST output
pixels are emitted.  Capturing on the coarse busy signal mixed those phases with
the convolution output phase and misaligned the trace-to-pixel mapping.  The
corrected trigger is:
\[
\texttt{tio\_trigger} = \texttt{out\_ce} \land (|\texttt{out\_we}|) \land
(\texttt{out\_index} < 784).
\]

\begin{table}[H]
\centering
\caption{Corrected MNIST hardware capture metadata.}
\begin{tabular}{ll}
\toprule
Field & Value \\
\midrule
Trigger source & First output write \\
Images captured & 500 \\
Kernels per image & 9 \\
Input shape & $28 \times 28 \times 1$ MNIST \\
Kernel shape & $3 \times 3$ \\
FPGA clock & 5 MHz \\
ADC sample rate & 105 MS/s \\
Samples per FPGA cycle & 21 \\
Samples per trace & 19,200 \\
Averaging & 20 captures per kernel trace \\
Cycle-zero offset & 62 cycles \\
\bottomrule
\end{tabular}
\end{table}

\section{Output Summary}
The corrected hardware capture and extraction pipeline was evaluated twice:
once with the paper-style transformation pipeline but no RC curve fitting, and
once with RC curve fitting enabled.  The table includes the recognition targets
reported by the paper for context.

\begin{table}[H]
\centering
\caption{MNIST reconstruction and recognition results.}
\begin{tabular}{lccc}
\toprule
Metric & No curve fit & Curve fit & Paper target \\
\midrule
Original-image recognition & 0.975 & 0.975 & -- \\
Background pixel accuracy & 0.586 & 0.635 & -- \\
Background recognition & 0.140 & 0.085 & 0.816 \\
Template mean pixel distance & 31.077 & 31.777 & -- \\
Template recognition & 0.130 & 0.130 & 0.898 \\
\bottomrule
\end{tabular}
\end{table}

Curve fitting improves the background pixel metric from 0.586 to 0.635, but it
does not improve the final attack objective.  Background recognition drops from
14.0\% to 8.5\%, and template recognition remains at 13.0\%.  This means the
curve-fit stage is not producing features that preserve digit-level structure
well enough for recognition.

\begin{figure}[H]
\centering
\includegraphics[width=0.92\linewidth]{figs/mnist_metric_summary.png}
\caption{Recognition and pixel-level metrics for the corrected MNIST hardware
run compared with the paper's reported recognition targets.}
\end{figure}

\section{Trace Transformation Pipeline}
The implemented processing path follows the paper's Section 5 structure as
closely as the CW-Lite data permits:

\begin{enumerate}
  \item Capture analog power traces from the CW305 while the BNN convolution
  core runs on hardware.
  \item Restore the trace baseline using the inverse high-pass approximation.
  \item Apply a low-pass filter.  The nominal paper cutoff is 60 MHz, but the
  CW-Lite's 105 MS/s sampling rate makes the usable cutoff lower than a 2.5 GHz
  oscilloscope setup.
  \item Align cycle windows using a Pearson-correlation cycle template.
  \item Sum or curve-fit each cycle to obtain one power feature per FPGA cycle.
  \item Apply the 62-cycle output offset and remap the extended $30 \times 30$
  raster into the $28 \times 28$ MNIST output grid.
\end{enumerate}

\begin{figure}[H]
\centering
\includegraphics[width=\linewidth]{figs/mnist_trace_transformations.png}
\caption{Actual hardware trace from image 0300 and kernel 0 as it moves through
raw capture, baseline restoration, low-pass filtering, cycle alignment,
per-cycle extraction, and output-grid remapping.}
\end{figure}

\section{Recovered Images}
The recovered images explain why recognition is poor.  The originals are clean
MNIST digits, but both reconstruction modes mostly produce row-structured or
sparse artifacts.  The template output often collapses to weak partial strokes
or biased shapes; it does not consistently reconstruct the original digit.

\begin{figure}[H]
\centering
\includegraphics[width=0.92\linewidth]{figs/mnist_recovery_examples.png}
\caption{Original MNIST digits and recovered images from the corrected hardware
captures.  The template column is independently scaled per image so weak
template artifacts remain visible.}
\end{figure}

\section{Why Performance Is Bad}
The failure is not because the MNIST classifier is broken.  The original-image
recognition rate is 97.5\%, and the hardware functional check passed with all
784 outputs bit-exact.  The failure is in the measured side-channel signal and
the resulting reconstruction features.

First, the measurement setup is much weaker than the paper setup.  The paper
uses a high-bandwidth oscilloscope workflow, while this run uses CW-Lite analog
capture at 105 MS/s.  At a 5 MHz FPGA clock this gives only about 21 samples per
cycle.  That is enough to segment cycles, but it is not equivalent to a 2.5 GHz
oscilloscope waveform.  The paper's 60 MHz low-pass step is also not faithfully
available here because the CW-Lite Nyquist limit is 52.5 MHz.

Second, the CW305 Artix-7 power measurement contains a large amount of unrelated
activity: clock network power, static FPGA power, I/O effects, board power
distribution effects, and non-convolution control activity.  The useful dynamic
power caused by a single MNIST convolution output bit is small compared with
that background.  The row bands visible in the remapped power image indicate
that board-level or schedule-level structure dominates over pixel-level leakage.

Third, the current HLS target is functionally correct but not identical to the
paper's target implementation.  The present core includes a fanout-based leakage
amplifier to make the signal more visible.  Even with that amplifier, the
recovered images do not reach the paper's recognition quality, which suggests
the acquisition path is the limiting factor.

Fourth, the curve-fitting transform is not automatically helpful at CW-Lite
resolution.  It was intended for denser analog waveforms.  In this run it raises
the background pixel score but damages background recognition, which means it is
changing the statistics in a way that does not preserve the digit shape.

\section{Review}
\textbf{What is correct now.}
The run is MNIST-aligned.  The HLS C-simulation passed, the CW305 functional
check passed, and the output trigger now starts at the first useful write rather
than the broad busy window.  The trace metadata records trigger source, sample
rate, FPGA clock, samples per cycle, averaging, and the 62-cycle remap offset.
The attack pipeline is reproducible from saved trace files and saved evaluation
directories.

\textbf{What remains risky.}
The implementation is aligned with the paper at the algorithmic-processing
level, but the measurement instrument and FPGA board are not paper-equivalent.
The CW-Lite bandwidth and ADC rate are the biggest mismatch.  The Artix-7 target
and CW305 power path are also likely noisier and more diluted than the smaller
FPGA setup in the paper.  The fanout leakage amplifier helps observability but
also means the target is not a strict reproduction of the paper's hardware.

\section{Problems}
\begin{enumerate}
  \item The old busy-window captures should be treated as invalid for MNIST
  reconstruction because they were not aligned to the first useful output cycle.
  \item The corrected captures are real hardware captures, but their signal to
  noise ratio is too low for high-quality digit recovery.
  \item The CW-Lite sampling rate prevents a faithful reproduction of the
  paper's 2.5 GHz oscilloscope processing path.
  \item The CW305 board-level power rail mixes the small BNN leakage with large
  unrelated FPGA and board activity.
  \item The template method is statistically unstable on this data and tends to
  produce sparse, biased reconstructions.
  \item The curve-fit stage is computationally expensive and, on this capture
  setup, does not improve recognition.
\end{enumerate}

\section{Recommended Next Steps}
For a stronger reproduction of the paper, the best next move is measurement
improvement rather than more software tuning.  Use a higher-bandwidth capture
instrument, a better-isolated power rail, an EM probe close to the active FPGA
region, or a smaller target implementation with less unrelated dynamic power.
If staying with CW-Lite, the next useful experiments are SNR/TVLA-style leakage
checks per cycle, frequency-domain inspection, and target modifications that
isolate only the convolution datapath during the capture window.

\section{Conclusion}
The MNIST hardware path is now functionally correct and actually running on the
CW305 hardware.  The bad recognition numbers are real and expected from the
current measurement quality.  The capture trigger and data alignment are fixed,
but the recovered images show that the trace features are dominated by
measurement noise and board-level structure rather than clean MNIST pixel
leakage.  The experiment is therefore aligned enough to diagnose the problem,
but not yet strong enough to reproduce the paper's recognition rates.

\end{document}
